\section{Background}
\label{sec:background}

\subsection{BFS-Based Partitioning in the Multilevel Literature}

Breadth-first search partitioning appears as an internal component of
several multilevel graph partitioning frameworks.
\citet{karypis1998} describe BFS-based initial partitioning in Section~4
of the METIS paper: after the graph has been coarsened to a few hundred
vertices, a BFS from two seeds produces an initial bisection that is
immediately refined by Kernighan-Lin boundary swaps.
The BFS step's output quality is deliberately not optimised---it exists
only to give KL a feasible starting point.

The use of BFS as a \emph{standalone} partitioner---without subsequent
KL refinement---has not been studied as a redistricting method.
This paper fills that gap by evaluating BFS growth as a complete
Structure-layer algorithm.

\subsection{k-Farthest Seed Selection}

The quality of a BFS partition is determined primarily by seed placement.
Seeds that are close together produce uneven districts (one large, one
small); seeds at opposite ends of the graph produce more balanced initial
boundaries.

The \emph{k-farthest} strategy selects seeds greedily by BFS distance:
\begin{enumerate}
\item Seed~0 is selected by population-weighted sampling (denser areas
  first, reflecting the fact that high-population areas are more likely
  to require splitting and thus benefit from being starting points).
\item Each subsequent seed is the tract with the maximum minimum BFS
  distance to any already-selected seed.
\end{enumerate}
This strategy is a BFS analogue of the farthest-point initialisation used
in $k$-means++ for Euclidean spaces~\citep{cormen2009}.
It ensures that seeds are maximally spread across the graph topology rather
than clustering in densely connected regions.

For the redistricting bisection case ($k=2$), the k-farthest strategy
reduces to: place seed~0 by weighted sampling, then place seed~1 at the
BFS-farthest tract from seed~0.
This is exactly the strategy that maximises the initial boundary length
between the two growing regions.

\subsection{The Population-Deficit Priority Function}

Once seeds are placed, the growth order determines the final district shapes.
\bfsg{} uses a \emph{minimum-priority queue} keyed on current district
population:
\[
  \text{priority}(\text{district } d) = P_d^{\mathrm{current}},
\]
where $P_d^{\mathrm{current}}$ is the total population already assigned
to district $d$.
The tract with the lowest priority key is expanded next; ties are broken
by tract index.

This priority function fills the \emph{emptiest} district first.
It is motivated by the following balance argument: if district $d$ has
population far below the ideal, its neighbours should be assigned to $d$
before they are assigned to a more populated district.
Formally, under the deficit-priority function, the final imbalance is
bounded by the maximum single-tract population, because each growth step
reduces the maximum deficit by at least the added tract's population.

An alternative priority function---$|\text{ideal} - P_d^{\mathrm{current}}|$
(absolute deficit)---would produce the same ordering when a district is
below ideal, but would deprioritise districts that have exceeded their
ideal population, potentially allowing overfull districts to grow further.
The current-population priority avoids this pathology by always filling
the smallest district.

\paragraph{Comparison with CVD.}
\cvd{} assigns each tract to the district whose centroid (geographic center
of mass) is closest, iterating until centroid positions stabilise~\citep{du1999}.
The centroidal Voronoi equilibrium minimises within-district sum of squared
geographic distances---a proxy for compactness.

\bfsg{} assigns tracts by BFS adjacency order, with population deficit
breaking ties.
It does not compute centroids and does not iterate.
The single-pass nature means that \bfsg{} cannot recover from early
suboptimal assignments, but it also means that \bfsg{} runs in $O(n \log n)$
rather than $O(n \times k \times \text{iter})$.

Table~\ref{tab:comparison-algorithms} summarises the relationship between
\bfsg{}, \cvd{}, and \metis{} along the key dimensions.

\begin{table}[t]
\centering
\caption{Comparison of structure-layer algorithms.}
\label{tab:comparison-algorithms}
\begin{tabular}{lcccc}
\toprule
Property & \bfsg{} & \cvd{} & \metis{} \\
\midrule
Metric optimised & None (greedy) & Geographic distance & Edge cut \\
Data needed & Adjacency only & Adjacency + centroids & Adjacency \\
Iterations & 1 pass & $\sim$20 iters & Multilevel KL \\
Runtime & $O(n \log n)$ & $O(n \times k \times 20)$ & $O(n \log n)$ \\
Contiguity guarantee & Yes (BFS) & No (post-process) & No (post-process) \\
Compactness & Low (floor) & Medium & Medium-high \\
Hyperparameters & 0 & 1 (iter count) & 1 (ufactor) \\
\bottomrule
\end{tabular}
\end{table}

\subsection{Structure-Layer Context in the B-Series Platform}

The B-series redistricting platform uses a three-layer compositor:
\emph{Structure} (how each bisection node is solved),
\emph{WeightsOverride} (what edge weights encode), and
\emph{SeedCompositor} (how many seeds are tried).
\bfsg{} is a Structure-layer algorithm.

Because \bfsg{} is a Structure-layer algorithm, it composes cleanly with
all SeedCompositor strategies.
Running \be{} with $T=50$ seeds, each using \bfsg{} growth, selects the
best of 50 BFS growth outcomes.
Running \cs{} with \bfsg{} sweeps forward through seeds until $T=600$
consecutive seeds produce no improvement.
In both cases, the SeedCompositor provides the multi-seed search
quality improvement that \bfsg{} alone cannot.

Existing structure-layer algorithms all modify \emph{what is optimised}
at each node:
\geosec{}~\citep{b8paper} minimises normalised edge cut;
\areasec{} adds area-swing constraints;
\ar{}~\citep{b11paper} links the bisection tree to Huntington-Hill
apportionment.
\bfsg{} is the unique structure-layer algorithm that optimises nothing:
it is the zero-optimisation baseline against which all others are measured.
